\section{习题}

\subsection{基础过关题}

\begin{enumerate}

\item \textbf{（单项选择题）} 对同一组记录分别执行直接插入排序和折半插入排序，二者都采用本章给出的代码（升序）。与直接插入排序相比，折半插入排序
\begin{center}
\begin{tabularx}{0.96\linewidth}{@{}X@{}}
(A) 比较次数与后移次数都减少 \\[4pt]
(B) 只减少比较次数，后移次数不变 \\[4pt]
(C) 只减少后移次数，比较次数不变 \\[4pt]
(D) 总时间复杂度降为 $O(n\log n)$
\end{tabularx}
\end{center}

\ifshowsolutions\par\noindent\textbf{解析：} 折半插入只改变“找插入位置”的方式，不改变“腾位置”的方式。

插入 $A[i]$ 时，折半查找的比较次数不超过 $\lceil\log_2(i+1)\rceil$，而直接插入最坏要比较 $i$ 次，因此 $n$ 趟累计的比较次数由 $\Theta(n^2)$ 降到 $\Theta(n\log n)$。但位置找到以后，仍要把 $A[low..i-1]$ 中的元素逐个后移一格；后移的元素恰好是“前缀中大于 $temp$ 的元素”，这个集合与直接插入排序完全相同，最后的落位写法也一样。所以移动次数一点没变，仍为 $\Theta(n^2)$。总时间由移动主导，还是 $O(n^2)$，故 (D) 错。

顺带注意：折半查找在 $A[mid]==temp$ 时令 $low=mid+1$，把相等记录插在其后，所以折半插入排序仍是稳定排序。

\boxed{\textbf{答案：}(B)}\else \answerarea \fi

\item \textbf{（单项选择题）} 对序列 $\{5,1,4,2,8,3\}$ 执行一趟冒泡排序：从前向后依次比较相邻元素，逆序即交换，扫描到末尾为止，暂不启用“本趟无交换则提前结束”的优化。第一趟结束后的序列是
\begin{center}
\begin{tabularx}{0.96\linewidth}{@{}XX@{}}
(A) $\{1,4,2,5,3,8\}$ & (B) $\{1,2,4,5,3,8\}$ \\[4pt]
(C) $\{1,4,2,3,5,8\}$ & (D) $\{1,4,2,5,8,3\}$
\end{tabularx}
\end{center}

\ifshowsolutions\par\noindent\textbf{解析：} 思路：一趟冒泡只保证把当前未排序部分的最大元素送到末尾，中途元素只能靠相邻交换逐步挪动，不会“跳”成完全有序。

按 $j=0,1,\ldots,4$ 依次比较：
\[
\begin{aligned}
\{5,1,4,2,8,3\} &\xrightarrow{5>1} \{1,5,4,2,8,3\}
\xrightarrow{5>4} \{1,4,5,2,8,3\}\\
&\xrightarrow{5>2} \{1,4,2,5,8,3\}
\xrightarrow{5<8} \{1,4,2,5,8,3\}\\
&\xrightarrow{8>3} \{1,4,2,5,3,8\}.
\end{aligned}
\]
共比较 $n-1=5$ 次，交换 4 次；最大值 8 落到最后一个位置，但前面的 $4$ 与 $2$ 仍未就位，这正是“一趟”与“有序”的区别。

\boxed{\textbf{答案：}(A)}\else \answerarea \fi

\item \textbf{（单项选择题）} 对序列 $\{46,25,68,13,57,32,79,21\}$，按本章“取首元素为枢轴、先从右端找小、再从左端找大、交替填坑”的划分算法执行一趟划分，结果为
\begin{center}
\begin{tabularx}{0.96\linewidth}{@{}XX@{}}
(A) $\{21,25,32,13,46,57,79,68\}$ & (B) $\{21,25,32,13,46,57,68,79\}$ \\[4pt]
(C) $\{13,25,32,21,46,57,79,68\}$ & (D) $\{21,25,13,32,46,57,79,68\}$
\end{tabularx}
\end{center}

\ifshowsolutions\par\noindent\textbf{解析：} 思路：枢轴 46 先被“挖出”，它在数组首位的格子成为第一个坑；此后右指针找小于枢轴的元素填左坑，左指针找大于枢轴的元素填右坑，两指针相遇处的坑放回枢轴。

下标 $0\sim 7$，初值 $low=0$、$high=7$，$pivot=46$：
\begin{itemize}
\item $A[7]=21<46$，填入 $A[0]$：$\{21,25,68,13,57,32,79,21\}$，坑移到下标 7；
\item $A[0]=21\leqslant 46$、$A[1]=25\leqslant 46$，$low$ 移到 2；$A[2]=68>46$，填入 $A[7]$：$\{21,25,68,13,57,32,79,68\}$，坑移到下标 2；
\item $A[7]=68\geqslant 46$、$A[6]=79\geqslant 46$，$high$ 移到 5；$A[5]=32<46$，填入 $A[2]$：$\{21,25,32,13,57,32,79,68\}$，坑移到下标 5；
\item $A[2]=32\leqslant 46$、$A[3]=13\leqslant 46$，$low$ 移到 4；$A[4]=57>46$，填入 $A[5]$：$\{21,25,32,13,57,57,79,68\}$，坑移到下标 4；
\item $A[5]=57\geqslant 46$，$high$ 减到 4，与 $low$ 相遇，$A[4]=46$ 落位。
\end{itemize}
结果为 $\{21,25,32,13,46,57,79,68\}$：枢轴在下标 4，左段全都 $\leqslant 46$，右段全都 $\geqslant 46$。注意此时左右两段内部并未排好，所以 (B) 把右段写成有序是错的。

\boxed{\textbf{答案：}(A)}\else \answerarea \fi

\item \textbf{（单项选择题）} 已知大根堆按层序存放在数组 $A[0..5]$ 中，序列为 $\{20,15,18,8,9,12\}$。删除堆顶元素后，将堆尾元素移到堆顶并向下调整，得到的堆是
\begin{center}
\begin{tabularx}{0.96\linewidth}{@{}XX@{}}
(A) $\{18,15,12,8,9\}$ & (B) $\{15,12,18,8,9\}$ \\[4pt]
(C) $\{18,15,9,8,12\}$ & (D) $\{18,12,15,8,9\}$
\end{tabularx}
\end{center}

\ifshowsolutions\par\noindent\textbf{解析：} 思路：删除堆顶后堆的规模变为 5。把最后一个元素 $12$ 放到根上，此时左右两棵子树仍是堆，只有根可能违反堆序，向下调整即可。每一步都要先选出左右孩子中的较大者，再拿当前结点与它比较。

原堆 $\{20,15,18,8,9,12\}$，删去 20，把 12 移到 $A[0]$：$\{12,15,18,8,9\}$。
根的两个孩子是 $A[1]=15$、$A[2]=18$，较大者为 18；$12<18$，交换：$\{18,15,12,8,9\}$。
结点 12 落在下标 2，其孩子下标 5、6 已越界，调整结束。

得到的序列满足 $18\geqslant 15,12$、$15\geqslant 8,9$，是大根堆。这里容易错选 (D)：如果先与左孩子 15 比较并交换，就会破坏“与较大孩子交换”的规则。

\boxed{\textbf{答案：}(A)}\else \answerarea \fi

\item \textbf{（填空题）} 对序列 $\{9,1,5,8,3,7,4,6,2\}$ 采用希尔排序，增量依次取 $4,2,1$（$d_1=\lfloor 9/2\rfloor=4$，此后每次取前一个增量的一半并向下取整）。第一趟（$d=4$）结束后的序列为 \underline{\qquad\qquad\qquad\qquad}。

\ifshowsolutions\par\noindent\textbf{解析：} 思路：增量为 $d$ 时，下标相差 $d$ 的元素属于同一个子序列，即第 $i$ 组为 $i,\,i+d,\,i+2d,\ldots$。每组各自做直接插入排序，其余位置的元素不动。

$d=4$ 时共 4 组：
\[
\{9,3,2\}\ \text{（下标 }0,4,8\text{）}\to\{2,3,9\},\qquad
\{1,7\}\ \text{（下标 }1,5\text{）}\to\{1,7\},
\]
\[
\{5,4\}\ \text{（下标 }2,6\text{）}\to\{4,5\},\qquad
\{8,6\}\ \text{（下标 }3,7\text{）}\to\{6,8\}.
\]
按原下标填回：下标 $0,1,2,3,4,5,6,7,8$ 依次为 $2,1,4,6,3,7,5,8,9$。

注意这一趟并没有让整个序列接近有序（前面仍是 $2,1$），希尔排序的威力来自后续 $d=2$、$d=1$ 的逐层整理。另外，相隔 $d$ 的相等关键字可能被分进不同的子序列，各子序列分别排序后它们的相对次序可能改变，所以希尔排序不稳定。

\boxed{\text{答案：}\{2,1,4,6,3,7,5,8,9\}}\else \answerarea \fi

\item \textbf{（填空题）} 将两个长度分别为 $4$ 和 $5$ 的有序表归并为一个有序表，最坏情况下需要比较 \underline{\qquad} 次；对 $n=1000$ 个记录做二路归并排序（自底向上，段长依次为 $1,2,4,\ldots$），共需归并 \underline{\qquad} 趟。

\ifshowsolutions\par\noindent\textbf{解析：} 思路：归并时每比较一次，至少有一个元素被取走；只有当某一段已经取空，另一段的剩余元素才能直接照抄、不再比较。最坏情况是“最后一个元素才决出胜负”，于是前 $m+n-1$ 个元素各比较一次，最后一个直接复制，故最多
\[
m+n-1=4+5-1=8
\]
次。最少比较次数是 $\min(m,n)=4$ 次（短表整体先被取空）。

趟数：每趟把段长翻倍，段长达到不小于 $n$ 时排序结束，故趟数为 $\lceil\log_2 n\rceil$。由
\[
2^9=512<1000\leqslant 1024=2^{10}
\]
得 $\lceil\log_2 1000\rceil=10$，即 10 趟。这个趟数与初始序列无关，是归并排序“任何输入都是 $O(n\log n)$”的直接原因。

\boxed{\textbf{答案：}\text{8 次；10 趟}}\else \answerarea \fi

\end{enumerate}

\subsection{真题风格与综合训练}

\begin{enumerate}[start=7]

\item \textbf{（真题风格·单项选择题）} 对序列 $\{329,457,657,839,436,720,355\}$ 按最低位优先（LSD）做十进制基数排序，第一趟按个位分配到 10 个队列后依次收集，所得序列是
\begin{center}
\begin{tabularx}{0.96\linewidth}{@{}XX@{}}
(A) $\{720,355,436,457,657,329,839\}$ & (B) $\{720,355,436,457,657,839,329\}$ \\[4pt]
(C) $\{720,355,436,657,457,329,839\}$ & (D) $\{329,355,436,457,657,720,839\}$
\end{tabularx}
\end{center}

\ifshowsolutions\par\noindent\textbf{解析：} 思路：LSD 第一趟只看个位，元素进桶的次序就是原序列次序，收集时按桶号 $0\to9$、桶内先进先出。

各元素的个位：$329\to9$、$457\to7$、$657\to7$、$839\to9$、$436\to6$、$720\to0$、$355\to5$。于是
\[
\text{桶 }0:\{720\},\ \text{桶 }5:\{355\},\ \text{桶 }6:\{436\},\
\text{桶 }7:\{457,657\},\ \text{桶 }9:\{329,839\}.
\]
依次收集得 $\{720,355,436,457,657,329,839\}$。

这里有两个易错点。第一，桶内必须保持原来的先后次序：457 原本在 657 之前，收集后仍在前，这正是“每一趟分配收集都稳定”的要求，否则前一趟的结果会被破坏，最终得不到有序序列。第二，一趟结束后序列只是“按个位有序”，整体并不有序，(D) 是完全排好的结果，属于概念混淆。

\boxed{\textbf{答案：}(A)}\else \answerarea \fi

\item \textbf{（真题风格·单项选择题）} 某文件经内部排序后得到 100 个初始归并段，采用 4 路平衡归并，并用败者树选出当前最小记录。建立败者树后，每输出一个记录平均需要的比较次数与总归并趟数分别为
\begin{center}
\begin{tabularx}{0.96\linewidth}{@{}XX@{}}
(A) 约 2 次，3 趟 & (B) 约 2 次，4 趟 \\[4pt]
(C) 约 4 次，4 趟 & (D) 约 100 次，25 趟
\end{tabularx}
\end{center}

\ifshowsolutions\par\noindent\textbf{解析：} 思路：趟数只由“每趟段数缩小为 $1/k$”决定，即 $k^t\geqslant m$；比较次数只由败者树的树高决定，与段数 $m$ 无关。

\textbf{趟数。} 4 路归并每趟把段数变为约原来的 $1/4$：
\[
4^3=64<100,\qquad 4^4=256\geqslant 100,
\]
故需 $t=\lceil\log_4 100\rceil=4$ 趟。若误用 $4^3=64$ 就判定 3 趟，是把“至少 64 段”当成了“覆盖 100 段”。

\textbf{比较次数。} 败者树有 $k=4$ 个叶结点，树高为 $\lceil\log_2 k\rceil=2$。输出一个记录后，只需从它所在的叶结点向上重赛到根，每层与父结点中保存的“失败者”比较一次，共 2 次；建树的 $O(k)$ 次比较只在开始时做一次。若不用败者树而每趟线性扫描 4 个段的首元素，每输出一个记录要比较 3 次，且 $k$ 越大差距越大——这正是败者树的价值。

\boxed{\textbf{答案：}(B)}\else \answerarea \fi

\item \textbf{（真题风格·综合题）} 设关键字序列 $S=\{52,31,74,18,63,45,87,26\}$，下标从 $0$ 开始，所有排序均按升序进行。约定：
\begin{itemize}
\item 直接插入排序采用本章代码，仅当 $A[j]>temp$ 时才后移；
\item 简单选择排序每趟在未排序部分中选最小元素，与未排序部分的第一个元素交换；
\item 二路归并排序采用自底向上方式，段长依次为 $1,2,4,8$，归并时相等元素先取左段；
\item 快速排序以区间首元素为枢轴，按本章“右找小、左找大、交替填坑”的划分算法。
\end{itemize}
回答：
\begin{enumerate}
\item[(1)] 直接插入排序执行完 $i=1,2,3$ 三趟后的序列；
\item[(2)] 简单选择排序执行完前两趟后的序列；
\item[(3)] 二路归并排序执行完第一趟、第二趟后的序列；
\item[(4)] 以 52 为枢轴执行一趟划分后的序列；
\item[(5)] 上述四种算法中哪些稳定、哪些不稳定，并各给出判断理由。
\end{enumerate}

\ifshowsolutions\par\noindent\textbf{解析：} 各小题都从初始序列 $S$ 出发独立执行，不接续上一小题的结果。

\textbf{(1) 直接插入排序。} 前 $i$ 个元素始终有序。
$i=1$：$temp=31$，$52>31$ 后移，得 $\{31,52,74,18,63,45,87,26\}$；
$i=2$：$temp=74$，$52<74$，不动；
$i=3$：$temp=18$，依次与 $74,52,31$ 比较并后移，得
\[
\{18,31,52,74,63,45,87,26\}.
\]

\textbf{(2) 简单选择排序。} 第 1 趟在全部 8 个元素中选最小者 18（下标 3），与 $A[0]=52$ 交换，得 $\{18,31,74,52,63,45,87,26\}$；第 2 趟在 $A[1..7]$ 中选最小者 26（下标 7），与 $A[1]=31$ 交换，得
\[
\{18,26,74,52,63,45,87,31\}.
\]

\textbf{(3) 二路归并排序。} 第一趟把相邻单元素段两两归并：
\[
(52,31)\to31,52;\quad(74,18)\to18,74;\quad(63,45)\to45,63;\quad(87,26)\to26,87,
\]
得 $\{31,52,18,74,45,63,26,87\}$。第二趟把相邻的长度 2 的段两两归并：
\[
[31,52]\cup[18,74]\to18,31,52,74;\qquad[45,63]\cup[26,87]\to26,45,63,87,
\]
得 $\{18,31,52,74,26,45,63,87\}$。（第三趟即可得到全局有序 $\{18,26,31,45,52,63,74,87\}$。）

\textbf{(4) 快速排序一趟划分。} $pivot=52$，$low=0$，$high=7$：
右端 $26<52$，填入 $A[0]$，坑移到 7；左端 $26,31\leqslant52$，$low$ 到 2，$A[2]=74>52$ 填入 $A[7]$，坑移到 2；右端 $74,87\geqslant52$，$high$ 到 5，$A[5]=45<52$ 填入 $A[2]$，坑移到 5；左端 $45,18\leqslant52$，$low$ 到 4，$A[4]=63>52$ 填入 $A[5]$，坑移到 4；右端 $A[5]=63\geqslant52$，$high$ 减到 4 与 $low$ 相遇，$A[4]=52$ 落位。得
\[
\{26,31,45,18,52,63,87,74\}.
\]

\textbf{(5) 稳定性。} 直接插入排序只在 $A[j]>temp$ 时后移，相等记录不会被越过，稳定；归并时相等元素先取左段，稳定。简单选择排序不稳定：取 $S'=\{2,2',1\}$，第 1 趟选最小元素 1 与第一个 2 交换，得 $\{1,2',2\}$，两个相等关键字的次序颠倒。快速排序不稳定：同样取 $S'=\{2,2',1\}$，以首元素 2 为枢轴划分，右端 1 填入首位，两个 2 中靠后的那个经枢轴落位后反而排到了前面，得 $\{1,2',2\}$。

\fbox{\parbox{0.93\linewidth}{\textbf{答案：}\\
(1) $\{18,31,52,74,63,45,87,26\}$；\\
(2) $\{18,26,74,52,63,45,87,31\}$；\\
(3) 第一趟 $\{31,52,18,74,45,63,26,87\}$，第二趟 $\{18,31,52,74,26,45,63,87\}$；\\
(4) $\{26,31,45,18,52,63,87,74\}$；\\
(5) 稳定：直接插入排序、二路归并排序；不稳定：简单选择排序、快速排序。}}\else \answerarea[10] \fi

\item \textbf{（真题风格·综合题）} 设计一个算法：从含 $n$ 个整数的数组 $A[0..n-1]$ 中找出最小的 $k$ 个元素（$1\leqslant k\leqslant n$），存入 $B[0..k-1]$，不要求 $B$ 内部有序。要求时间复杂度为 $O(n\log k)$，并回答：
\begin{enumerate}
\item[(1)] 说明算法思想与正确性依据；
\item[(2)] 写出 C 语言代码（可调用向下调整函数 \texttt{SiftDown}）；
\item[(3)] 分析时间、空间复杂度，并与“先整体排序再取前 $k$ 个”“每趟线性扫描选最小”作比较。
\end{enumerate}

\ifshowsolutions\par\noindent\textbf{解析：} \textbf{(1) 算法思想。} 维护一个大小为 $k$ 的候选集合，它始终保存“已扫描元素中最小的 $k$ 个”。为了 $O(1)$ 判断新元素是否值得进入集合，需要随时取出集合中的最大者；为了替换后快速恢复，需要 $O(\log k)$ 的调整——这正好是大根堆：堆顶就是这 $k$ 个候选中的最大者。

先把 $A[0..k-1]$ 建成大根堆。此后扫描 $A[k..n-1]$：若 $A[i]\geqslant$ 堆顶，说明已有 $k$ 个元素不大于它，它不可能进入前 $k$ 小，直接跳过；若 $A[i]<$ 堆顶，就用它替换堆顶并向下调整。扫描结束时，堆中恰为最小的 $k$ 个元素。

\textbf{正确性（循环不变式）。} 当扫描到 $A[i]$（$i\geqslant k$）时，$B$ 中恰为 $A[0..i-1]$ 中最小的 $k$ 个元素，且 $B[0]$ 是其中最大者。
初始时 $B$ 为前 $k$ 个元素，建堆后不变式成立。归纳步：若 $A[i]\geqslant B[0]$，则至少已有 $k$ 个元素不大于 $A[i]$，它不属于最小的 $k$ 个；若 $A[i]<B[0]$，则 $A[i]$ 必然替换当前集合中的最大者，替换后集合仍是 $k$ 个最小元素，再下沉恢复堆序即保持 $B[0]$ 为最大。

\textbf{(2) 代码。}
\begin{lstlisting}[language=C]
// 大根堆向下调整：把 H[i] 下沉，堆的大小为 m
void SiftDown(int H[], int i, int m) {
    int temp = H[i];
    int j = 2 * i + 1;                 // 左孩子
    while (j < m) {
        if (j + 1 < m && H[j] < H[j + 1])
            j++;                       // 取较大的孩子
        if (temp >= H[j]) break;       // 已满足堆序
        H[i] = H[j];                   // 大孩子上移
        i = j;
        j = 2 * i + 1;
    }
    H[i] = temp;
}

void MinK(int A[], int n, int k, int B[]) {
    int i;
    for (i = 0; i < k; i++) B[i] = A[i];   // 前 k 个先入堆
    for (i = k / 2 - 1; i >= 0; i--)       // 自底向上建大根堆
        SiftDown(B, i, k);
    for (i = k; i < n; i++) {
        if (A[i] < B[0]) {                 // 比堆中最大者还小才替换
            B[0] = A[i];
            SiftDown(B, 0, k);             // 恢复大根堆
        }
    }
}
\end{lstlisting}

\textbf{(3) 复杂度与比较。} 建堆 $O(k)$；其余 $n-k$ 个元素中，每个最多引起一次替换，每次下沉沿树高走 $O(\log k)$，故总时间为
\[
O(k)+O\bigl((n-k)\log k\bigr)=O(n\log k).
\]
额外空间只有 \texttt{temp}、\texttt{i}、\texttt{j} 等常数个变量（$B$ 是输出，不计入辅助空间），即 $O(1)$。

与另外两种做法比较：整体排序再取前 $k$ 个是 $O(n\log n)$，当 $k\ll n$ 时多做无用功；每趟线性扫描选最小、共做 $k$ 趟是 $O(nk)$，$k$ 稍大就不如本算法。本算法把“候选集合”限制在 $k$ 个，比较和调整都只在这个小堆内进行，所以得到了 $O(n\log k)$。当 $k=n$ 时它退化为堆排序，复杂度 $O(n\log n)$，与理论一致。

\fbox{\parbox{0.93\linewidth}{\textbf{答案：}用大小为 $k$ 的大根堆保存当前最小的 $k$ 个元素：先建堆，再依次用更小的元素替换堆顶并下沉。时间 $O(n\log k)$，辅助空间 $O(1)$（$B$ 为输出）。}}\else \answerarea[10] \fi

\end{enumerate}
